% Part: proof-theory
% Chapter: normalization
% Section: reductions

\documentclass[../../../include/open-logic-section]{subfiles}

\begin{document}

\olfileid{pt}{nor}{red}

\olsection{হ্রাস-রূপান্তর}

\Log{N1i}-র !!a{proof}-এ কোনো
কর্তনের দৈর্ঘ্য~$1$ হলে সেই খণ্ডে
মাত্র একটি !!{formula}-সংঘটন থাকে।
ওই !!{formula} একই সঙ্গে
!!a{introduction} বিধির সিদ্ধান্ত
এবং !!a{elimination} বিধির প্রধান
পূর্বধারণা। স্বাভাবিকীকরণের প্রমাণে
এমন কর্তন \emph{হ্রাস} করতে ওই
!!{introduction}/!!{elimination}
জোড়ায় শেষ হওয়া উপ-!!{proof}-এর
জায়গায় নতুন !!{proof} বসাই,
যেখানে বিধিজোড়াটি বাদ গেছে।
এতে নতুন কর্তন তৈরি হতে পারে;
তবে তাদের মাত্রা হ্রাসকৃত
কর্তনের মাত্রার চেয়ে কম।
ফলে নতুন !!{proof}-এর
কর্তনদৈর্ঘ্য কমে (দৈর্ঘ্য~$1$-এর
সর্বাধিক কর্তনটি বাদ গেছে)
অথবা !!{proof}-এর
কর্তনমাত্রা কমে (যদি সেটিই
একমাত্র সর্বাধিক কর্তন হয়ে থাকে)।

যেমন, সংশ্লিষ্ট বিধি
\Intro{\lif} ও \Elim{\lif}
হলে, $!A \ident !B \lif !C$
হলে এবং উপ-!!{proof}~$\delta_1$
নিচের মতো শেষ হলে:
\[
\AxiomC{$\Discharge{!B}{x}$}
\RightLabel{$\delta_2$}
\DeduceC{$!C$}
\DischargeRule{\Intro{\lif}}{x}
\UnaryInfC{$!B \lif !C$}
\AxiomC{}
\RightLabel{$\delta_3$}
\DeduceC{$!B$}
\RightLabel{\Elim{\lif}}
\BinaryInfC{$!C$}
\DisplayProof
\]
এই কর্তন সরাতে $\delta$-র
উপপ্রমাণ~$\delta_1$-র
জায়গায় নিচেরটি বসাই:
\[
\AxiomC{}
\RightLabel{$\delta_3$}
\DeduceC{$!B$}
\RightLabel{$\Subst{\delta_2}{\delta_3}{!B^x}$}
\DeduceC{$!C$}
\DisplayProof
\]
$\Subst{\delta_2}{\delta_3}{!B^x}$
বলতে সেই !!{proof} বোঝাই,
যা~$\delta_2$-তে
$x$-চিহ্নযুক্ত প্রতিটি
খোলা অনুমান~$!B$-র জায়গায়
পুরো !!{proof}~$\delta_3$
বসিয়ে পাওয়া যায়; কলম-স্থাপনের আগের
অব্যাহতি-চিহ্ন পরিবর্তনের রীতিও মেনে নিই।
এই !!{proof}-এ মূল প্রমাণের প্রতিস্থাপিত
$x$-চিহ্নযুক্ত $!B$-রূপের অনুমানগুলি আর থাকে না;
তবে বসানো প্রমাণের নিজস্ব খোলা অনুমান থাকতে পারে।
সেই অনুমানও~$x$-চিহ্নযুক্ত হতে পারে, কারণ একই চিহ্নে
দুটি ভিন্ন !!{formula}s
অনুমান হিসেবে আসে না।
তাই নতুন !!{proof}-এর খোলা
অনুমানগুলি~$\delta_1$-র খোলা অনুমানগুলির উপসেট;
প্রতিস্থাপনযোগ্য খোলা অনুমান না থাকলে অপ্রয়োজনীয় শাখা বাদ যায়;
$\delta_2$-র যে বিধিগুলি
অনুমান অবমুক্ত করে, তারা
$\Subst{\delta_2}{\delta_3}{!B^x}$-এও
একই অনুমান অবমুক্ত করে;
আর $\delta_3$-র অবমুক্তিকারী
বিধিগুলির একাধিক অনুলিপি
হতে পারে, যেগুলি
$\Subst{\delta_2}{\delta_3}{!B^x}$-এ
সংশ্লিষ্ট অনুমান অবমুক্ত করে।
% BN-SRC-880 TARGET-CORRECTION-BEGIN
\par\smallskip\noindent\textnormal{\small\textbf{পাঠ-সংশোধন (BN-SRC-880–882):}
কলম-স্থাপন মূল উপপ্রমাণের খোলা অনুমানগুলিতেই হয়। বসানো
প্রমাণের নিজস্ব খোলা অনুমান থাকতে পারে; তাই একই চিহ্নের কোনো
খোলা অনুমানই আর নেই বলা যায় না। শূন্য অব্যাহতির ক্ষেত্রে
অপ্রয়োজনীয় শাখা বাদ যায়, ফলে খোলা অনুমানের সমতা নয়,
উপসেট হওয়াই সাধারণভাবে নিশ্চিত।}
% BN-SRC-880 TARGET-CORRECTION-END
% BN-SRC-881: inserted proof may still have its own open assumption with the substituted label.
% BN-SRC-882: vacuous discharge can discard the entire unused premise; open assumptions form a subset.


$\delta_2$ বা $\delta_3$-র সম্পূর্ণ
ভিতরে, অথবা $\delta$-তে
$\delta_1$-র সম্পূর্ণ বাইরে থাকা
কর্তন অপরিবর্তিত: তাতে একই
!!{formula}s ও একই বিধি থাকে।
তাই $\delta$-র সংশ্লিষ্ট কর্তনের
মাত্রা ও দৈর্ঘ্যও একই থাকে।
দৈর্ঘ্য~$1$-এর একটি সর্বাধিক
কর্তন সরিয়েছি; ফলে হয়
কর্তনদৈর্ঘ্য কমেছে, নয়
কর্তনমাত্রা কমেছে (যদি
$\delta$-র কর্তনদৈর্ঘ্য~$1$
ছিল এবং এটিই একমাত্র
সর্বাধিক কর্তন ছিল)।
% BN-SRC-578 TARGET-CORRECTION-BEGIN
\par\smallskip\noindent\textnormal{\small\textbf{উৎস-সংশোধন (BN-SRC-578):} এই →-হ্রাস উদাহরণে উপপ্রমাণ কেবল $\delta_2$ ও $\delta_3$; উৎসের অপরিবর্তিত-কর্তনের তালিকায় অনির্ধারিত $\delta_4$ বাদ দেওয়া হয়েছে।}
% BN-SRC-578 TARGET-CORRECTION-END

তবে দুটি ঘটনা ঘটতে পারে:
\begin{enumerate}
  \item সব $!B^x$ অনুমানের জায়গায়
  $\delta_3$ বসালে
  $\delta_3$-র সব কর্তনের
  অনুলিপি বেড়ে যায়। তাদের
  কোনোটি সর্বাধিক কর্তন হলে
  নতুন !!{proof}-এর
  কর্তনদৈর্ঘ্য পুরোনো
  !!{proof}-এর চেয়ে বাড়তে পারে।
  তা না ঘটার নিশ্চয়তা দরকার।
  তাই শুধু \emph{সবচেয়ে ডানের}
  সর্বাধিক কর্তনে হ্রাস-রূপান্তর
  প্রয়োগ করি: যে কর্তনের
  ডান দিকে~$\delta$-র কোনো
  শাখায় আর সর্বাধিক কর্তন
  নেই। $\delta_3$-তে কর্তন
  থাকলে তা এখানে হ্রাসকৃত
  কর্তনের ডান দিকে পড়ত;
  তখন আমরা সবচেয়ে ডানের
  কর্তন নিয়ে কাজ করতাম না।
  সব সময় সবচেয়ে ডানের
  কর্তন নিলে হয় !!{proof}-এর
  কর্তনদৈর্ঘ্য কমে, নয়
  কর্তনমাত্রাই কমে।
  \item $\delta_2$-র কোনো
  অনুমান~$!B^x$ যদি
  !!a{elimination} বিধির
  প্রধান পূর্বধারণা হয়,
  আর $\delta_3$-র শেষ বিধি
  \Elim{\lor}, \Elim{\lexists}
  বা !!a{introduction} হয়,
  তাহলে ওই অনুমানে
  $\delta_3$-কে $\delta_2$-তে
  কলম-স্থাপন করে নতুন
  কর্তন তৈরি করি।
  $\delta_2$-তে যা শুধু
  অনুমান ছিল, এখন তা
  দৈর্ঘ্য~$1$-এর কর্তন
  (একটি !!a{introduction}
  বিধির সিদ্ধান্ত এবং
  !!a{elimination} বিধির
  প্রধান পূর্বধারণা),
  অথবা কর্তনখণ্ডের শেষ
  !!{formula}-সংঘটন
  (!!a{elimination} বিধির
  প্রধান পূর্বধারণা এবং
  \Elim{\lor} বা
  \Elim{\lexists}-এর সিদ্ধান্ত)।
  $!B^x$ যদি \Elim{\lor}
  বা \Elim{\lexists}-এর
  গৌণ পূর্বধারণা হয়ে থাকে,
  তবে তা আগে থেকেই একটি
  খণ্ডের প্রথম !!{formula}-সংঘটন,
  সম্ভবত কর্তনখণ্ডেরও।
  নতুন !!{proof}-এ খণ্ডটি
  আরও দীর্ঘ হতে পারে।
  কিন্তু নতুন কর্তন বা পুরোনো
  কর্তনখণ্ড গড়া !!{formula}
  হলো~$!B$, এবং
  $\depth{!B} < \depth{!B \lif !C}$।
  অতএব কর্তন যোগ হওয়া বা
  দৈর্ঘ্য বাড়ার সম্ভাবনা
  থাকলেও এগুলির কোনোটিই
  সর্বাধিক কর্তন নয়।
  ফলে হয় কর্তনমাত্রা কমেছে,
  নয় একই মাত্রার কর্তন
  থেকে গেলে কর্তনদৈর্ঘ্য কমেছে।
  % BN-SRC-579 TARGET-CORRECTION-BEGIN
  \par\smallskip\noindent\textnormal{\small\textbf{উৎস-সংশোধন (BN-SRC-579):} এই একই অনুমান আগে $!B^x$ হলেও গৌণ-পূর্বধারণার বাক্যে উৎসে ভুলে $B^x$ লেখা; সূত্রচিহ্নসহ $!B^x$ রাখা হয়েছে।}
  % BN-SRC-579 TARGET-CORRECTION-END
\end{enumerate}

দৈর্ঘ্য~$1$-এর কর্তন হ্রাসের
প্রক্রিয়াকে \emph{হ্রাস-রূপান্তর}
(বা \emph{ঘুরপথ-রূপান্তর}) বলে।
কয়েকটি বিধিজোড়ার ছক
\olref{tab:reduction-conversions}-এ
দেওয়া হয়েছে। (অন্যগুলি
অনুশীলনের জন্য রাখা।)

\begin{table}
\begin{multline*}
\bottomAlignProof
\AxiomC{$\Discharge{!B}{x}$}
\RightLabel{$\delta_2$}
\shortDeduce
\DeduceC{$!C$}
\DischargeRule{\Intro{\lif}}{x}
\UnaryInfC{$!B \lif !C$}
\AxiomC{}
\RightLabel{$\delta_3$}
\shortDeduce
\DeduceC{$!B$}
\RightLabel{\Elim{\lif}}
\BinaryInfC{$!C$}
\DisplayProof
\quad \longrightarrow\quad
\shoveright{\bottomAlignProof
\AxiomC{}
\RightLabel{$\delta_3$}
\shortDeduce
\DeduceC{$!B$}
\RightLabel{$\Subst{\delta_2}{\delta_3}{!B^x}$}
\shortDeduce
\DeduceC{$!C$}
\DisplayProof}
\\[4ex]
\shoveleft{\bottomAlignProof
\AxiomC{}
\RightLabel{$\delta_2$}
\shortDeduce
\DeduceC{$!B$}
\AxiomC{}
\RightLabel{$\delta_3$}
\shortDeduce
\DeduceC{$!C$}
\RightLabel{\Intro{\land}}
\BinaryInfC{$!B \land !C$}
\RightLabel{\Elim{\land}[1]}
\UnaryInfC{$!B$}
\DisplayProof
\quad \longrightarrow\quad
\bottomAlignProof
\AxiomC{}
\RightLabel{$\delta_2$}
\shortDeduce
\DeduceC{$!B$}
\DisplayProof}\qquad
\shoveright{\bottomAlignProof
\AxiomC{}
\RightLabel{$\delta_2$}
\shortDeduce
\DeduceC{$!B$}
\AxiomC{}
\RightLabel{$\delta_3$}
\shortDeduce
\DeduceC{$!C$}
\RightLabel{\Intro{\land}}
\BinaryInfC{$!B \land !C$}
\RightLabel{\Elim{\land}[2]}
\UnaryInfC{$!C$}
\DisplayProof
\quad \longrightarrow\quad
\bottomAlignProof
\AxiomC{}
\RightLabel{$\delta_3$}
\shortDeduce
\DeduceC{$!C$}
\DisplayProof}
\\[4ex]
\bottomAlignProof
\AxiomC{}
\RightLabel{$\delta_2$}
\shortDeduce
\DeduceC{$!B$}
\RightLabel{\Intro{\lor}[1]}
\UnaryInfC{$!B \lor !C$}
\AxiomC{$\Discharge{!B}{x}$}
\RightLabel{$\delta_3$}
\shortDeduce
\DeduceC{$!D$}
\AxiomC{$\Discharge{!C}{y}$}
\RightLabel{$\delta_4$}
\shortDeduce
\DeduceC{$!D$}
\DischargeRule{\Elim{\lor}}{x\,y}
\TrinaryInfC{$!D$}
\DisplayProof
\quad \longrightarrow\quad
\shoveright{\bottomAlignProof
\AxiomC{}
\RightLabel{$\delta_2$}
\shortDeduce
\DeduceC{$!B$}
\RightLabel{$\Subst{\delta_3}{\delta_2}{!B^x}$}
\shortDeduce
\DeduceC{$!D$}
\DisplayProof}
\\[2ex]
\shoveleft{\bottomAlignProof
\AxiomC{}
\RightLabel{$\delta_2$}
\shortDeduce
\DeduceC{$!B(c)$}
\RightLabel{\Intro{\lforall}}
\UnaryInfC{$\lforall[x][!B(x)]$}
\RightLabel{\Elim{\lforall}}
\UnaryInfC{$!B(t)$}
\DisplayProof
\quad \longrightarrow\quad
\bottomAlignProof
\AxiomC{}
\RightLabel{$\Subst{\delta_2}{t}{c}$}
\shortDeduce
\DeduceC{$!B(t)$}
\DisplayProof}
\\[4ex]
\shoveright{\bottomAlignProof
\AxiomC{}
\RightLabel{$\delta_2$}
\shortDeduce
\DeduceC{$!B(t)$}
\RightLabel{\Intro{\lexists}}
\UnaryInfC{$\lexists[x][!B(x)]$}
\AxiomC{$\Discharge{!B(c)}{x}$}
\RightLabel{$\delta_3$}
\shortDeduce
\DeduceC{$!D$}
\DischargeRule{\Elim{\lexists}}{x}
\BinaryInfC{$!D$}
\DisplayProof
\quad \longrightarrow\quad
\bottomAlignProof
\AxiomC{}
\RightLabel{$\delta_2$}
\shortDeduce
\DeduceC{$!B(t)$}
\RightLabel{$\Subst{\Subst{\delta_3}{t}{c}}{\delta_2}{!B(t)^x}$}
\shortDeduce
\DeduceC{$!D$}
\DisplayProof}
\end{multline*}
\caption{\Log{N1i}-এর কয়েকটি হ্রাস-রূপান্তর।}
\ollabel{tab:reduction-conversions}
\end{table}

\begin{prob}
\Intro{\lnot}-এর পরে
\Elim{\lnot} থাকলে
হ্রাস-রূপান্তরটি দিন।
\end{prob}

আরও একটি ক্ষেত্র বিচার করতে হবে।
কর্তনের সংজ্ঞায় দৈর্ঘ্য~$1$-এর
এমন ক্ষেত্রও আছে, যেখানে
$!A$ একদিকে \FalseInt{}-এর
সিদ্ধান্ত, অন্যদিকে
!!a{elimination} বিধির
প্রধান পূর্বধারণা।
$!A$ যখন !!a{introduction}
বিধির সিদ্ধান্ত, তার মতোই
এ ক্ষেত্র সামলাই। যেমন,
নিচেরটি বদলে
\[
\bottomAlignProof
\AxiomC{}
\RightLabel{$\delta_2$}
\DeduceC{$\lfalse$}
\UnaryInfC{$!B \lif !C$}
\AxiomC{}
\RightLabel{$\delta_3$}
\DeduceC{$!B$}
\RightLabel{\Elim{\lif}}
\BinaryInfC{$!C$}
\DisplayProof
\quad\text{এর বদলে}\quad
\bottomAlignProof
\AxiomC{}
\RightLabel{$\delta_2$}
\DeduceC{$\lfalse$}
\UnaryInfC{$!C$}
\DisplayProof
\]
% BN-SRC-580: optional explicit FalseInt labels withdrawn; unlabelled displayed steps already determine the rule.
$!A \ident (!B \lor !C)$
অথবা $!A \ident \lexists[x][!B(x)]$
হলে একটু সতর্ক থাকতে হয়।
যেমন, নিচেরটি বদলে
\[
\bottomAlignProof
\AxiomC{}
\RightLabel{$\delta_2$}
\DeduceC{$\lfalse$}
\RightLabel{\FalseInt}
\UnaryInfC{$!B \lor !C$}
\AxiomC{$\Discharge{!B}{x}$}
\RightLabel{$\delta_3$}
\DeduceC{$!D$}
\AxiomC{$\Discharge{!C}{y}$}
\RightLabel{$\delta_4$}
\DeduceC{$!D$}
\DischargeRule{\Elim{\lor}}{x\,y}
\TrinaryInfC{$!D$}
\DisplayProof
\quad\text{এর বদলে}\quad
\bottomAlignProof
\AxiomC{}
\RightLabel{$\delta_2$}
\DeduceC{$\lfalse$}
\RightLabel{\FalseInt}
\UnaryInfC{$!D$}
\DisplayProof
\]
% BN-SRC-581 TARGET-CORRECTION-BEGIN
\par\smallskip\noindent\textnormal{\small\textbf{উৎস-সংশোধন (BN-SRC-581):} $\lor$-নিষ্কাশনের দুই ভিন্ন অনুমান উৎসে একই $x$ চিহ্ন পেয়েছে; N1i-তে ভিন্ন সূত্রে একই অনুমান-চিহ্ন না-থাকার শর্ত অনুযায়ী $!C$-শাখায় $y$ এবং নিষ্কাশনবিধিতে $x,y$ অবমুক্তি স্পষ্ট করা হয়েছে।}
% BN-SRC-581 TARGET-CORRECTION-END
তাহলে কর্তন-!!{formula}~$!D$-যুক্ত
নতুন কর্তন তৈরি হতে পারে;
কিন্তু $\depth{!D} < \depth{!B \lor !C}$
হবেই, এমন নিশ্চয়তা নেই!{}
তাই \Intro{\lor}/\Elim{\lor}
হ্রাসের মতোই এখানে নিচেরটি
বসাতে হবে:
\[
\bottomAlignProof
\AxiomC{}
\RightLabel{$\delta_2$}
\DeduceC{$\lfalse$}
\RightLabel{\FalseInt}
\UnaryInfC{$!B$}
\RightLabel{$\delta_3$}
\DeduceC{$!D$}
\DisplayProof
\quad\text{অথবা}\quad
\bottomAlignProof
\AxiomC{}
\RightLabel{$\delta_2$}
\DeduceC{$\lfalse$}
\RightLabel{\FalseInt}
\UnaryInfC{$!C$}
\RightLabel{$\delta_4$}
\DeduceC{$!D$}
\DisplayProof
\]
% BN-NORM-173: display connectors remain localized; explanatory reader annotation withdrawn.

\end{document}
